\documentclass[10pt, aspectratio=169]{beamer}
\usepackage{../beamer_theme_408}

\title{408考研零基础 --- 数据结构}
\subtitle{1-11 二叉树的遍历}
\author{}
\date{}


\begin{document}


\begin{frame}{本节目标}
    \begin{enumerate}
        \item 理解先序、中序、后序遍历的递归定义与访问顺序
        \item 掌握非递归遍历算法的实现
        \item 掌握层序遍历的队列实现，能根据遍历序列还原二叉树
    \end{enumerate}
\end{frame}

\begin{frame}{递归遍历}
    \textbf{先序遍历（根左右）}
    \[
    \text{根} \to \text{左子树} \to \text{右子树}
    \]
    \textbf{中序遍历（左根右）}
    \[
    \text{左子树} \to \text{根} \to \text{右子树}
    \]
    \textbf{后序遍历（左右根）}
    \[
    \text{左子树} \to \text{右子树} \to \text{根}
    \]
    \vspace{0.5em}
    \begin{block}{核心区别}
        根结点的访问时机决定遍历方式名称
    \end{block}
\end{frame}

\begin{frame}[fragile]{递归算法实现}
    \textbf{先序遍历}
    \begin{lstlisting}[language=C, basicstyle=\ttfamily\small]
void PreOrder(BiTree T) {
    if(T != NULL) {
        visit(T);
        PreOrder(T->lchild);
        PreOrder(T->rchild);
    }
}
    \end{lstlisting}
    \vspace{0.3em}
    \textbf{中序和后序仅调整 visit 位置}
    \vspace{0.3em}
    时间复杂度 $O(n)$，空间复杂度 $O(h)$
\end{frame}

\begin{frame}{非递归遍历 --- 先序}
    \textbf{使用栈模拟}
    \vspace{0.3em}
    \begin{enumerate}
        \item 根结点入栈
        \item 出栈 $p$，访问 $p$
        \item 右孩子入栈（先右后左，保证左先出）
        \item 左孩子入栈
        \item 重复直到栈空
    \end{enumerate}
\end{frame}

\begin{frame}{非递归遍历 --- 中序}
    \textbf{使用栈模拟}
    \vspace{0.3em}
    \begin{enumerate}
        \item $p$ 指向根
        \item 若 $p$ 非空，$p$ 入栈，$p = p \to \text{lchild}$
        \item 若 $p$ 为空且栈非空，出栈 $p$，访问 $p$
        \item $p = p \to \text{rchild}$
        \item 重复直到栈空且 $p$ 为空
    \end{enumerate}
\end{frame}

\begin{frame}{层序遍历}
    \textbf{使用队列}
    \vspace{0.3em}
    \begin{enumerate}
        \item 根结点入队
        \item 出队 $p$，访问 $p$
        \item 左孩子入队
        \item 右孩子入队
        \item 重复直到队空
    \end{enumerate}
    \vspace{0.5em}
    \textbf{时间复杂度} $O(n)$，\textbf{空间复杂度} $O(n)$
    \vspace{0.3em}
    \textbf{应用}：求最大宽度、判断完全二叉树
\end{frame}

\begin{frame}{由遍历序列还原二叉树}
    \textbf{先序 + 中序 $\implies$ 唯一二叉树}
    \vspace{0.3em}
    \textbf{原理}
    \begin{enumerate}
        \item 先序第一个为根
        \item 中序中根左边为左子树，右边为右子树
        \item 递归构建
    \end{enumerate}
    \vspace{0.3em}
    \textbf{例题}
    \begin{itemize}
        \item 先序：A B D E C F
        \item 中序：D B E A F C
    \end{itemize}
    根 A，左子树 (B D E) / (D B E)，右子树 (C F) / (F C)
\end{frame}

\begin{frame}{本节总结}
    \begin{enumerate}
        \item 先序（根左右）、中序（左根右）、后序（左右根）
        \item 递归遍历 $O(n)$ 时间，$O(h)$ 空间
        \item 非递归遍历用栈模拟递归
        \item 层序遍历用队列实现
        \item 先序/后序 + 中序可唯一还原二叉树
    \end{enumerate}
\end{frame}

\end{document}
